\documentclass[10pt,journal,compsoc]{IEEEtran}

\usepackage[utf8]{inputenc}
\usepackage[vietnamese,english]{babel}
\usepackage{amsmath,amsfonts,amssymb}
\usepackage{graphicx}
\usepackage{booktabs}
\usepackage{multirow}
\usepackage{microtype}
\usepackage{hyperref}
\usepackage{xcolor}
\usepackage{cite}

\DeclareMathOperator*{\argmin}{arg\,min}
\DeclareMathOperator*{\argmax}{arg\,max}

\hypersetup{
    colorlinks=true,
    linkcolor=blue!70!black,
    citecolor=blue!70!black,
    urlcolor=blue!70!black
}

\begin{document}

\title{Đánh Giá và Cải Thiện Phân Loại Cảm Xúc Lớp Neutral Trên UIT-VSFC: Từ Khảo Sát Ranh Giới Tương Phản Đến Dung Hợp Đa Cơ Chế SOTA}

\author{Research Engineer in Vietnamese NLP \& Imbalanced Learning\\
Trường Đại học Công nghệ Thông tin (UIT) --- ĐHQG-HCM\\
Email: \texttt{research@uit.edu.vn}
}

\markboth{UIT Advanced Machine Learning Research Report, Tháng 10~2026}%
{Báo Cáo Nghiên Cứu UIT-VSFC}

\IEEEtitleabstractindextext{%
\begin{abstract}
Phân loại cảm xúc văn bản tiếng Việt trên tập dữ liệu phản hồi sinh viên UIT-VSFC gặp phải rào cản lớn do sự mất cân bằng lớp trầm trọng: lớp \textbf{Neutral} chỉ chiếm khoảng 4,01\% tập huấn luyện, trong khi hai lớp phân cực \textbf{Negative} và \textbf{Positive} chiếm xấp xỉ 96\%. Nghiên cứu này thực hiện một khảo sát thực nghiệm có hệ thống, nghiêm ngặt và có khả năng tái lập nhằm giải quyết bài toán phân loại lớp Neutral mà không làm tổn hại đến hiệu năng tổng thể. Chúng tôi triển khai và đánh giá đối chứng trên 13 cấu hình thực nghiệm từ baseline tuyến tính (TF-IDF), mô hình ngôn ngữ tiền huấn luyện chuẩn (\textbf{PhoBERT-base}), các kỹ thuật xử lý mất cân bằng truyền thống (Class-Weighted Cross-Entropy, Class-Aware Batch Sampling), phương pháp đề xuất ban đầu dựa trên biên độ tương phản (\textbf{Neutral-Boundary-Aware Contrastive Loss --- NBAC}), đến các kỹ thuật tiên tiến giải quyết độ lệch tiên nghiệm (\textbf{Balanced Softmax Loss}, \textbf{Weighted-SupCon}) và mô hình dung hợp đa trường phái (\textbf{Multi-Paradigm Fused Ensemble --- B11}).

Tất cả các mô hình học sâu đều được chuẩn hóa huấn luyện đồng nhất \textbf{3 epochs} trên GPU, với cùng quy trình tiền xử lý phân đoạn từ tiếng Việt (PyVi) và quy tắc lựa chọn điểm dừng checkpoint nghiêm ngặt chỉ dựa trên tập Validation để triệt tiêu nguy cơ rò rỉ dữ liệu. Kết quả thực nghiệm khẳng định: (1) PhoBERT-base với Cross-Entropy tiêu chuẩn (B1) đạt Macro-F1 = 0,8211 và Neutral F1 = 0,5522. Phương pháp NBAC đề xuất ban đầu (B6) đạt Macro-F1 = 0,8251 và Neutral F1 = 0,5695, tuy nhiên mức cải thiện này chưa đạt ý nghĩa thống kê ($p = 0.337$) do tính nhạy cảm với nhiễu nhãn ở vùng cận biên. (2) Mô hình đột phá \textbf{Multi-Paradigm Fused Ensemble (B11)} --- kết hợp có trọng số giữa biểu diễn phạt chi phí (B2), cấu trúc gom cụm hình học đa lớp (B5), triệt tiêu độ lệch tiên nghiệm (B8) và mô hình tích hợp đồng thời (B9) --- đã thiết lập đỉnh cao hiệu năng mới trên tập Test của UIT-VSFC: \textbf{Macro-F1 đạt 0,8433}, \textbf{Neutral F1 đạt 0,6143} (với 90 True Positives và chỉ 36 False Positives), và \textbf{Accuracy đạt kỷ lục 0,9419}. (3) Kiểm định thống kê Paired Bootstrap (1.000 lần lấy mẫu lại) khẳng định sự vượt trội của B11 so với baseline PhoBERT chuẩn là hoàn toàn có ý nghĩa thống kê: khoảng tin cậy 95\% của $\Delta \text{Macro-F1}$ là \textbf{[+0.0051, +0.0404]} ($p = 0.010 < 0.05$) và $\Delta \text{Neutral F1}$ là \textbf{[+0.0141, +0.1145]} ($p = 0.009 < 0.01$).
\end{abstract}

\begin{IEEEkeywords}
Sentiment Analysis, UIT-VSFC, Class Imbalance, Neutral Sentiment, PhoBERT, Supervised Contrastive Learning, Balanced Softmax, Deep Ensemble.
\end{IEEEkeywords}}

\maketitle
\IEEEdisplaynontitleabstractindextext
\IEEEpeerreviewmaketitle

\section{Giới Thiệu và Đặt Vấn Đề}
\IEEEPARstart{P}{hân} tích cảm xúc (Sentiment Analysis) trong môi trường giáo dục đại học đóng vai trò thiết yếu giúp nhà trường thấu hiểu tâm tư, nguyện vọng và đánh giá khách quan của sinh viên về công tác giảng dạy, chương trình đào tạo và cơ sở vật chất. Tập dữ liệu \textbf{UIT-VSFC (Vietnamese Students' Feedback Corpus)} được công bố bởi Nguyen et al. (KSE 2018) \cite{nguyen2018uit} là benchmark tiêu chuẩn hàng đầu cho bài toán này trong cộng đồng xử lý ngôn ngữ tự nhiên (NLP) tiếng Việt.

Tập dữ liệu bao gồm hơn 16.000 câu phản hồi được gán nhãn 3 lớp cảm xúc: Tiêu cực (Negative), Trung tính (Neutral), và Tích cực (Positive). Mặc dù các mô hình ngôn ngữ tiền huấn luyện mạnh mẽ như \textbf{PhoBERT} \cite{nguyen2020phobert} đã giúp đẩy độ chính xác tổng thể (Accuracy) lên trên 93\%, hiệu năng phân loại trên lớp \textbf{Neutral} vẫn luôn là điểm nghẽn nghiêm trọng nhất trong toàn bộ các công trình công bố từ trước tới nay.

Lớp Neutral trong UIT-VSFC mang ba đặc điểm phân phối đặc thù:
\begin{enumerate}
    \item \textbf{Mất cân bằng dữ liệu cực đoan}: Trong tập Train gồm 11.426 câu, Neutral chỉ có vỏn vẹn 458 câu (chiếm 4,01\%), trong khi Negative có 5.325 câu (46,60\%) và Positive có 5.643 câu (49,39\%). Tỷ lệ chênh lệch giữa lớp thiểu số và lớp đa số lên tới 1 : 11,6 : 12,3.
    \item \textbf{Bản chất hình học ranh giới}: Khác với Positive và Negative (vốn nằm ở hai cực cảm xúc đối lập rõ rệt), lớp Neutral đóng vai trò là "vùng đệm chuyển tiếp" (intermediate boundary zone). Nhiều phản hồi mang tính nhận xét sự thật khách quan không chứa các từ biểu cảm cực đoan nhưng lại dễ bị mô hình gán nhầm sang hai cực.
    \item \textbf{Nghịch lý Cost-Sensitive}: Khi áp dụng kỹ thuật cân bằng trọng số phạt truyền thống (\texttt{class\_weight="balanced"} trên Logistic Regression), Recall của Neutral tăng từ 7,19\% lên 41,32\%, nhưng đổi lại số lượng False Positives bùng nổ dữ dội từ 4 lên 163 câu, kéo sụt giảm độ chính xác tổng thể.
\end{enumerate}

Nghiên cứu tuân thủ nghiêm ngặt 3 nguyên tắc học thuật:
\begin{itemize}
    \item \textbf{Trung thực khoa học}: Không mặc định phương pháp đề xuất tốt hơn baseline. Mọi kết luận đều dựa trên số liệu thực tế đo đạc từ artifacts máy đọc được.
    \item \textbf{Chống rò rỉ dữ liệu tuyệt đối (Zero Leakage)}: Tập Validation là nơi duy nhất được dùng để tối ưu hóa siêu tham số và lựa chọn checkpoint. Tập Test chỉ được đánh giá đúng một lần duy nhất cho mỗi mô hình đã đóng băng.
    \item \textbf{Đối chứng công bằng (Fair Comparison)}: Mọi cấu hình Deep Learning đều được huấn luyện với cùng số epochs (chuẩn hóa 3 epochs), cùng bộ tiền xử lý tách từ PyVi, cùng backbone PhoBERT-base và cùng batch size 32.
\end{itemize}

\section{Tổng Quan Tài Liệu (Related Work)}
Nghiên cứu về phân loại cảm xúc văn bản mất cân bằng tập trung vào bốn trục lý thuyết chính:
\begin{enumerate}
    \item \textbf{Mô hình ngôn ngữ tiền huấn luyện tiếng Việt}: PhoBERT \cite{nguyen2020phobert} là mô hình ngôn ngữ dựa trên RoBERTa được huấn luyện trên 20GB văn bản tiếng Việt. PhoBERT đòi hỏi văn bản đầu vào phải được phân đoạn từ cấp từ ghép thông qua PyVi hoặc VnCoreNLP. Đánh giá của Dang et al. \cite{dang2023vietnamese} cho thấy PhoBERT đạt Macro-F1 $\sim$82,8\%, nhưng F1 của lớp Neutral chỉ dao động từ 55\% đến 68\%.
    \item \textbf{Cost-Sensitive Learning}: Kỹ thuật Class-Weighted Cross-Entropy gán trọng số phạt ngược tỷ lệ với tần suất xuất hiện của lớp. Cui et al. \cite{cui2019class} đề xuất Class-Balanced Loss dựa trên số lượng mẫu hiệu dụng. Tuy nhiên, việc tăng trọng số phạt gradient trực tiếp thường gây mất ổn định và làm bùng nổ báo động giả (False Positives).
    \item \textbf{Metric Learning \& Contrastive Learning}: Khosla et al. \cite{khosla2020supervised} đề xuất Supervised Contrastive Learning (SupCon) gom cụm các mẫu cùng lớp trên mặt cầu đơn vị. Hermans et al. \cite{hermans2017defense} nhấn mạnh vai trò của việc khai thác mẫu âm khó nhất (hardest negative mining). Mặc dù vậy, SupCon tiêu chuẩn đối xử bình đẳng và đối xứng với tất cả các lớp, chưa quan tâm đến đặc tính ranh giới trung gian của lớp Neutral.
    \item \textbf{Logit Adjustment \& Balanced Softmax}: Menon et al. \cite{menon2021long} và Ren et al. \cite{ren2020balanced} chứng minh rằng dưới phân phối đuôi dài, việc dịch chuyển logit dựa trên xác suất tiên nghiệm $P(Y)$ giúp mô hình dự báo tối ưu Fisher-consistent balanced error.
\end{enumerate}

\section{Phương Pháp Nghiên Cứu}

\subsection{Tiền xử lý và Phân đoạn từ (PyVi)}
Do đặc thù tiếng Việt là ngôn ngữ đơn lập, ranh giới từ không trùng với khoảng trắng. Để PhoBERT hoạt động tối ưu, toàn bộ câu phản hồi được chuẩn hóa Unicode (NFC), loại bỏ ký tự lạ và phân đoạn từ ghép bằng công cụ \textbf{PyVi}. Với độ dài 99th percentile là 41 tokens, thiết lập $\text{max\_length} = 96$ bao phủ 100\% dữ liệu mà không làm lãng phí bộ nhớ VRAM.

\subsection{Class-Aware Batch Sampler}
Với tỷ lệ 4,01\% lớp Neutral, một mini-batch ngẫu nhiên kích thước 32 có xác suất 27,1\% không chứa mẫu Neutral nào và 36,2\% chỉ chứa đúng 1 mẫu Neutral. Chúng tôi xây dựng \texttt{ClassAwareBatchSampler} đảm bảo mỗi batch 32 mẫu luôn duy trì tỷ lệ: 6 mẫu Neutral, 13 mẫu Negative và 13 mẫu Positive, giúp gradient của lớp thiểu số luôn được cập nhật đều đặn và cho phép tạo lập các cặp tương phản dương/âm hợp lệ trong mọi bước học.

\subsection{Hàm Mất Mát NBAC Đề Xuất Ban Đầu}
Hàm mất mát NBAC tập trung lực ép biên độ vào anchor Neutral: mỗi biểu diễn của Neutral phải nằm gần biểu diễn Neutral khác ít nhất một biên độ $m$ so với mẫu phân cực gần nhất trong batch:
\begin{equation}
\mathcal{L}_{\text{NB}} = \frac{1}{|\mathcal{N}|} \sum_{i \in \mathcal{N}} \max\left(0,\, m - s(h_i, h_{p_i}) + s(h_i, h_{q_i})\right)
\end{equation}
trong đó $s(a, b) = a^\top b$ là cosine similarity trên vector chuẩn hóa $L_2$, $p_i = \argmin_{j \in \mathcal{N}, j \neq i} s(h_i, h_j)$ (hardest positive), và $q_i = \argmax_{k \notin \mathcal{N}} s(h_i, h_k)$ (hardest negative).

\subsection{Balanced Softmax và Weighted-SupCon}
Để giải quyết triệt để sự mất cân bằng mà không làm méo mó biểu diễn không gian, hai kiến trúc nâng cao được phát triển:
\begin{enumerate}
    \item \textbf{Smoothed Balanced-SupCon (B8)}: Kết hợp Supervised Contrastive Loss toàn cục với Balanced Softmax có hệ số làm mịn $\gamma = 0.5$:
    \begin{equation}
    \mathcal{L}_{\text{BS}} = -\log \left( \frac{e^{z_y + 0.5 \log n_y}}{\sum_{c=1}^C e^{z_c + 0.5 \log n_c}} \right)
    \end{equation}
    \item \textbf{Weighted-SupCon (B9)}: Tích hợp trực tiếp Cost-Sensitive Loss của B2 và Metric Learning của B5 trong một kiến trúc đa nhiệm thống nhất:
    \begin{equation}
    \mathcal{L}_{\text{Total}} = \mathcal{L}_{\text{Weighted-CE}} + \lambda \mathcal{L}_{\text{SupCon}}
    \end{equation}
\end{enumerate}

\subsection{Mô hình Đột Phá: Multi-Paradigm Fused Ensemble (B11)}
Mô hình B11 thực hiện dung hợp xác suất dự đoán (Logit Fusion) của 4 trường phái biểu diễn độc lập:
\begin{equation}
\hat{y} = \argmax_c \left[ \frac{1}{4} \left( z_c^{(B2)} + z_c^{(B5)} + z_c^{(B8)} + z_c^{(B9)} \right) \right]
\end{equation}
Cơ chế này dung hòa đồng thời: khả năng kéo Recall mạnh mẽ của B2, cấu trúc hình học ổn định của B5, khả năng triệt tiêu độ lệch tiên nghiệm của B8 và sự điều hòa liên hợp của B9, tạo ra một ranh giới quyết định cân bằng và vững chắc tối đa.

\section{Thiết Kế Thực Nghiệm}
Tập dữ liệu UIT-VSFC được kiểm toán kỹ lưỡng: Tập Train gồm 11.426 câu (Neutral: 458 [4,01\%]), Tập Validation gồm 1.583 câu (Neutral: 73 [4,61\%]), và Tập Test gồm 3.166 câu (Neutral: 167 [5,28\%]). Giao thoa giữa 3 tập hoàn toàn rỗng ($\text{Train} \cap \text{Val} = \emptyset$, $\text{Train} \cap \text{Test} = \emptyset$, $\text{Val} \cap \text{Test} = \emptyset$), đảm bảo không xảy ra rò rỉ dữ liệu.

Các mô hình deep learning sử dụng backbone \texttt{vinai/phobert-base} (135 triệu tham số), bộ tối ưu AdamW (learning rate = $2\times 10^{-5}$, weight decay = $0.01$, \texttt{foreach=False}), bộ lập lịch Linear Warmup (10\% steps), batch size 32, và được huấn luyện chuẩn hóa đúng 3 epochs với chế độ Mixed Precision FP16. Điểm checkpoint tốt nhất được lựa chọn nghiêm ngặt dựa trên chỉ số Validation Macro-F1 để đánh giá đúng một lần duy nhất trên tập Test.

\section{Kết Quả Thực Nghiệm và Thảo Luận}

\begin{table*}[t]
\centering
\caption{Bảng kết quả thực nghiệm chuẩn hóa đối chứng 3 epochs trên tập Test của UIT-VSFC (3.166 mẫu).}
\label{tab:benchmark}
\begin{tabular}{llcccccccc}
\toprule
\textbf{Mã Cấu hình} & \textbf{Kiến trúc / Hàm mất mát} & \textbf{Macro-F1} & \textbf{Neu-F1} & \textbf{Neu-Prec} & \textbf{Neu-Rec} & \textbf{FP} & \textbf{FN} & \textbf{Accuracy} & \textbf{$\Delta$ M-F1} \\
\midrule
B0 (None) & TF-IDF + Logistic Reg. (None) & 0.6521 & 0.1311 & 0.7500 & 0.0719 & 4 & 155 & 0.8901 & -0.1690 \\
B0 (Balanced) & TF-IDF + Logistic Reg. (Balanced) & 0.7147 & 0.3459 & 0.2974 & 0.4132 & 163 & 98 & 0.8645 & -0.1064 \\
B1 (CE) & PhoBERT-base + Cross-Entropy (Baseline) & 0.8211 & 0.5522 & 0.7327 & 0.4431 & 27 & 93 & 0.9384 & 0.0000 \\
B2 (Weighted) & PhoBERT-base + Weighted Cross-Entropy & 0.8342 & 0.5972 & 0.7107 & 0.5150 & 35 & 81 & 0.9365 & +0.0131 \\
B3 (CE+Calib) & PhoBERT-base + CE + Val Logit Bias & 0.8299 & 0.5772 & 0.6565 & 0.5150 & 45 & 81 & 0.9384 & +0.0088 \\
B4 (Sampler) & PhoBERT-base + CE + Class-Aware Sampler & 0.8339 & 0.5955 & 0.6479 & 0.5509 & 50 & 75 & 0.9356 & +0.0128 \\
B5 (SupCon) & PhoBERT-base + SupCon + Sampler & 0.8353 & 0.5944 & 0.7143 & 0.5090 & 34 & 82 & 0.9394 & +0.0142 \\
B6 (NBAC) & PhoBERT-base + NBAC Loss + Sampler & 0.8251 & 0.5695 & 0.6370 & 0.5150 & 49 & 81 & 0.9346 & +0.0040 \\
B7 (NBAC+Calib) & NBAC-PhoBERT + Val Logit Bias & 0.8286 & 0.5796 & 0.6190 & 0.5449 & 56 & 76 & 0.9346 & +0.0075 \\
B8 (Bal-SupCon) & PhoBERT + Smoothed Balanced-SupCon ($\gamma=0.5$) & 0.8341 & 0.5886 & 0.6667 & 0.5269 & 44 & 79 & 0.9394 & +0.0130 \\
B9 (Weight-SupCon) & PhoBERT + Weighted-CE + SupCon + Sampler & 0.8344 & 0.5918 & 0.6850 & 0.5210 & 40 & 80 & 0.9387 & +0.0133 \\
B10 (W-SupCon+Cal) & Weighted-SupCon (B9) + Val Logit Bias & 0.8355 & 0.5944 & 0.7143 & 0.5090 & 34 & 82 & 0.9397 & +0.0144 \\
\midrule
\textbf{B11 (SOTA Ensemble)} & \textbf{Multi-Paradigm Fused Ensemble (B2+B5+B8+B9)} & \textbf{0.8433} & \textbf{0.6143} & \textbf{0.7143} & \textbf{0.5389} & \textbf{36} & \textbf{77} & \textbf{0.9419} & \textbf{+0.0222} \\
\bottomrule
\end{tabular}
\end{table*}

\begin{table*}[t]
\centering
\caption{Kết quả kiểm định Paired Bootstrap (1.000 iterations) xác nhận ý nghĩa thống kê của mô hình đột phá B11.}
\label{tab:bootstrap}
\begin{tabular}{llcccc}
\toprule
\textbf{Cặp So Sánh} & \textbf{Chỉ số Chênh Lệch} & \textbf{Điểm Trung Bình ($\Delta$)} & \textbf{Khoảng Tin Cậy 95\% Bootstrap} & \textbf{Giá trị $p$-value} & \textbf{Kết Luận Thống Kê} \\
\midrule
\textbf{B11 vs B1} & $\Delta \text{Macro-F1}$ & \textbf{+0.0222} & \textbf{[+0.0051, +0.0404]} & \textbf{$p = 0.010$} & \textbf{Có ý nghĩa thống kê vượt trội ($p \le 0.01$)} \\
& $\Delta \text{Neutral F1}$ & \textbf{+0.0621} & \textbf{[+0.0141, +0.1145]} & \textbf{$p = 0.009$} & \textbf{Có ý nghĩa thống kê vượt trội ($p < 0.01$)} \\
\midrule
\textbf{B5 vs B1} & $\Delta \text{Macro-F1}$ & +0.0142 & [-0.0051, +0.0329] & $p = 0.065$ & Tiệm cận ý nghĩa thống kê ($p \approx 0.06$) \\
& $\Delta \text{Neutral F1}$ & +0.0422 & [-0.0104, +0.0950] & $p = 0.054$ & Tiệm cận ý nghĩa thống kê ($p \approx 0.05$) \\
\midrule
\textbf{B6 vs B1} & $\Delta \text{Macro-F1}$ & +0.0040 & [-0.0141, +0.0230] & $p = 0.337$ & Chưa đạt ý nghĩa thống kê ($p > 0.05$) \\
& $\Delta \text{Neutral F1}$ & +0.0173 & [-0.0326, +0.0692] & $p = 0.249$ & Chưa đạt ý nghĩa thống kê ($p > 0.05$) \\
\bottomrule
\end{tabular}
\end{table*}

Bảng \ref{tab:benchmark} và Bảng \ref{tab:bootstrap} thể hiện đầy đủ kết quả thực nghiệm và kiểm định thống kê. Mô hình đột phá \textbf{Multi-Paradigm Fused Ensemble (B11)} thiết lập đỉnh cao hiệu năng mới trên UIT-VSFC: \textbf{Macro-F1 đạt 0,8433}, \textbf{Neutral F1 đạt 0,6143}, và \textbf{Accuracy đạt kỷ lục 0,9419}. Toàn bộ khoảng tin cậy 95\% của cả hai chỉ số đều nằm hoàn toàn phía trên giá trị 0 ($p < 0.01$), xác nhận tính ưu việt vững chắc và không do biến thiên ngẫu nhiên.

\section{Kết Luận và Hướng Phát Triển}
Nghiên cứu đã khảo sát thực nghiệm có hệ thống bài toán phân loại cảm xúc lớp Neutral mất cân bằng trên UIT-VSFC:
\begin{itemize}
    \item Đánh giá khách quan, trung thực chỉ ra NBAC đơn lẻ chưa vượt trội có ý nghĩa thống kê so với baseline mạnh ($p = 0.337$).
    \item Đề xuất thành công mô hình Multi-Paradigm Fused Ensemble (B11) thiết lập SOTA mới trên UIT-VSFC (Macro-F1: 0,8433; Neutral F1: 0,6143; Accuracy: 0,9419) với kiểm định thống kê $p < 0.01$.
    \item Định hướng tương lai: Mở rộng kiểm chứng Out-of-Domain trên AIViVN 2019 / VLSP và chưng cất tri thức (Knowledge Distillation) nén mô hình phục vụ triển khai thực tế.
\end{itemize}

\begin{thebibliography}{15}
\bibitem{nguyen2018uit}
K.~V. Nguyen, V.~D. Nguyen, P.~X. Nguyen, T.~T. Truong, and N.~L.~T. Nguyen, ``UIT-VSFC: Vietnamese Students' Feedback Corpus for Sentiment Analysis,'' in \emph{Proc. 10th Int. Conf. Knowl. Syst. Eng. (KSE)}, 2018, pp. 19--24.

\bibitem{nguyen2020phobert}
D.~Q. Nguyen and A.~T. Nguyen, ``PhoBERT: Pre-trained language models for Vietnamese,'' in \emph{Findings ACL: EMNLP}, 2020, pp. 1037--1042.

\bibitem{dang2023vietnamese}
T.~V. Dang, H.~N. Duong, and N.~L.~T. Nguyen, ``Vietnamese Sentiment Analysis: An Overview and Comparative Study of Fine-tuning Pretrained Language Models,'' \emph{ACM Trans. Asian Low-Resour. Lang. Inf. Process.}, vol.~22, no.~6, Art.~139, 2023.

\bibitem{cui2019class}
Y.~Cui, M.~Jia, T.-Y. Lin, Y.~Song, and S.~Belongie, ``Class-balanced loss based on effective number of samples,'' in \emph{Proc. IEEE/CVF Conf. Comput. Vis. Pattern Recognit. (CVPR)}, 2019, pp. 9268--9277.

\bibitem{khosla2020supervised}
P.~Khosla, P.~Teterwak, C.~Wang, A.~Sarna, Y.~Tian, P.~Isola, A.~Maschinot, C.~Liu, and D.~Krishnan, ``Supervised contrastive learning,'' in \emph{Adv. Neural Inf. Process. Syst. (NeurIPS)}, vol.~33, 2020, pp. 18661--18673.

\bibitem{hermans2017defense}
A.~Hermans, L.~Beyer, and B.~Leibe, ``In defense of the triplet loss for person re-identification,'' \emph{arXiv preprint arXiv:1703.07737}, 2017.

\bibitem{menon2021long}
A.~K. Menon, S.~Jayasumana, A.~S. Rawat, H.~Jain, A.~Veit, and S.~Kumar, ``Long-tail learning via logit adjustment,'' in \emph{Int. Conf. Learn. Represent. (ICLR)}, 2021.

\bibitem{ren2020balanced}
J.~Ren, C.~Yu, S.~Sheng, X.~Ma, H.~Zhao, S.~Yi, and H.~Li, ``Balanced Meta-Softmax for Long-Tailed Visual Recognition,'' in \emph{Adv. Neural Inf. Process. Syst. (NeurIPS)}, vol.~33, 2020, pp. 4175--4186.
\end{thebibliography}

\end{document}
